\begin{conclusion}[采样对角权的结果式--判别式乘积公式：\texorpdfstring{$\prod_m K(b_m,b_m)$}{prod K(bm,bm)} 与 \texorpdfstring{$\prod_j K(x_j,x_j)$}{prod K(xj,xj)} 的闭式]\label{con:xi-terminal-wronskian-diagonal-sampling-resultant-discriminant-product}
沿用结论 \ref{con:xi-terminal-wronskian-kernel-discrete-mercer-parseval} 的记号，
并记 \(x_1,\dots,x_\kappa\) 为 \(y\) 的互异实根、\(b_0,\dots,b_\kappa\) 为 \(y_{+}\) 的互异实根。
则成立严格乘积恒等式
\[
\boxed{\;
\prod_{m=0}^{\kappa}K(b_m,b_m)
=\pi^{-(\kappa+1)}\operatorname{Res}(A,y_{+})
=\pi^{-(\kappa+1)}(-1)^{\frac{\kappa(\kappa+1)}{2}}\operatorname{Res}(y,y_{+})\,\mathrm{Disc}(y_{+})\;},
\]
以及
\[
\boxed{\;
\prod_{j=1}^{\kappa}K(x_j,x_j)
=\pi^{-\kappa}\operatorname{Res}(A,y)
=\pi^{-\kappa}(-1)^{\frac{\kappa(\kappa+1)}{2}}\operatorname{Res}(y,y_{+})\,\mathrm{Disc}(y)\;},
\]
其中
\(
\mathrm{Disc}(y)=\prod_{1\le i<j\le \kappa}(x_i-x_j)^2
\)、
\(
\mathrm{Disc}(y_{+})=\prod_{0\le i<j\le \kappa}(b_i-b_j)^2
\)。
\end{conclusion}
\begin{proof}[证明]
由 \(K(x,x)=A(x)/\pi\)（结论 \ref{con:xi-terminal-wronskian-kernel-discrete-mercer-parseval}）得
\[
\prod_{m=0}^{\kappa}K(b_m,b_m)=\pi^{-(\kappa+1)}\prod_{m=0}^{\kappa}A(b_m)
=\pi^{-(\kappa+1)}\operatorname{Res}(A,y_{+}).
\]
并由结论 \ref{con:xi-terminal-wronskian-resultant-discriminant-factorization} 得
\(
\operatorname{Res}(A,y_{+})
=(-1)^{\frac{\kappa(\kappa+1)}{2}}\operatorname{Res}(y,y_{+})\,\mathrm{Disc}(y_{+})
\)，代回即得第一条恒等式。
\par
同理，
\[
\prod_{j=1}^{\kappa}K(x_j,x_j)=\pi^{-\kappa}\prod_{j=1}^{\kappa}A(x_j)
=\pi^{-\kappa}\operatorname{Res}(A,y).
\]
又因 \(y(x_j)=0\)，代入 Wronskian 恒等式 \(A=yy_{+}'-y'y_{+}\) 得
\(
A(x_j)=-y'(x_j)\,y_{+}(x_j)
\)。
因此
\[
\operatorname{Res}(A,y)
=\prod_{j=1}^{\kappa}A(x_j)
=(-1)^{\kappa}\Bigl(\prod_{j=1}^{\kappa}y'(x_j)\Bigr)\Bigl(\prod_{j=1}^{\kappa}y_{+}(x_j)\Bigr)
=(-1)^{\kappa}\Bigl(\prod_{j=1}^{\kappa}y'(x_j)\Bigr)\operatorname{Res}(y,y_{+}).
\]
对首一多项式 \(y\) 有
\(
\prod_{j=1}^{\kappa}y'(x_j)=(-1)^{\frac{\kappa(\kappa-1)}{2}}\mathrm{Disc}(y)
\)，
合并符号 \( (-1)^{\kappa+\frac{\kappa(\kappa-1)}{2}}=(-1)^{\frac{\kappa(\kappa+1)}{2}} \) 即得
\(
\operatorname{Res}(A,y)
=(-1)^{\frac{\kappa(\kappa+1)}{2}}\operatorname{Res}(y,y_{+})\,\mathrm{Disc}(y)
\)。
代回即得第二条恒等式。证毕。
\end{proof}

\begin{conclusion}[纯 Lagrange 形式的核恒等式：\texorpdfstring{$K$}{K} 的显式采样展开与双变量代数恒等式]\label{con:xi-terminal-wronskian-kernel-lagrange-sampling-identity}
沿用结论 \ref{con:xi-terminal-wronskian-kernel-discrete-mercer-parseval} 的设定，并令 \(b_0,\dots,b_\kappa\) 为 \(y_{+}\) 的互异实根。
定义 Lagrange 基多项式
\[
L_m(z):=\frac{y_{+}(z)}{(z-b_m)\,y_{+}'(b_m)}\qquad (m=0,\dots,\kappa).
\]
则对一切 \(x,y\in\RR\) 成立显式核展开
\[
\boxed{\;
K(x,y)=\sum_{m=0}^{\kappa}K(b_m,b_m)\,L_m(x)\,L_m(y)
=\sum_{m=0}^{\kappa}\frac{A(b_m)}{\pi}\,L_m(x)\,L_m(y)\;}
\]
并等价地成立纯代数双变量恒等式
\[
\boxed{\;
\frac{y_{+}(x)y(y)-y(x)y_{+}(y)}{x-y}
=y_{+}(x)\,y_{+}(y)\sum_{m=0}^{\kappa}\frac{y(b_m)}{(x-b_m)(y-b_m)\,y_{+}'(b_m)}\; }.
\]
\end{conclusion}
\begin{proof}[证明]
由 \(y_{+}(b_m)=0\) 得
\[
K(x,b_m)=\frac{y_{+}(x)y(b_m)}{\pi(x-b_m)},
\qquad
K(b_m,b_m)=\frac{A(b_m)}{\pi}=\frac{y(b_m)\,y_{+}'(b_m)}{\pi},
\]
从而 \(K(x,b_m)=K(b_m,b_m)\,L_m(x)\)。
将其代入结论 \ref{con:xi-terminal-wronskian-kernel-discrete-mercer-parseval} 的离散 Mercer 分解即得
\(
K(x,y)=\sum_m K(b_m,b_m)L_m(x)L_m(y)
\)。
双变量形式由 \(K\) 的定义与 \(K(b_m,b_m)=y(b_m)y_{+}'(b_m)/\pi\) 整理得到。证毕。
\end{proof}

\begin{conclusion}[顶系数抽取公式：\texorpdfstring{$[z^{\kappa}]$}{[z^k]} 的 Lagrange--导数权表示]\label{con:xi-terminal-wronskian-top-coefficient-extraction}
沿用结论 \ref{con:xi-terminal-wronskian-kernel-lagrange-sampling-identity} 的记号。
令 \(q\in\CC[z]\) 且 \(\deg q\le \kappa\)。
则成立顶系数抽取恒等式
\[
\boxed{\;[z^{\kappa}]\,q(z)=\sum_{m=0}^{\kappa}\frac{q(b_m)}{y_{+}'(b_m)}\; }.
\]
特别地，
\[
\sum_{m=0}^{\kappa}\frac{1}{y_{+}'(b_m)}=0,\qquad
\sum_{m=0}^{\kappa}\frac{b_m}{y_{+}'(b_m)}=0,\qquad
\sum_{m=0}^{\kappa}\frac{y(b_m)}{y_{+}'(b_m)}=1.
\]
\end{conclusion}
\begin{proof}[证明]
对 \(\deg q\le\kappa\) 作 Lagrange 插值：
\[
q(z)=\sum_{m=0}^{\kappa} q(b_m)\,L_m(z)
=\sum_{m=0}^{\kappa}q(b_m)\,\frac{y_{+}(z)}{(z-b_m)\,y_{+}'(b_m)}.
\]
注意 \(y_{+}(z)/(z-b_m)\) 为首一次数 \(\kappa\) 多项式，因此 \(L_m\) 的 \(z^\kappa\) 系数为 \(1/y_{+}'(b_m)\)。
比较两边的 \(z^\kappa\) 系数即得主恒等式。
对 \(q(z)=1\)、\(q(z)=z\)、\(q(z)=y(z)\) 分别应用，并用 \([z^\kappa]1=0\)、\([z^\kappa]z=0\)、\([z^\kappa]y(z)=1\)，即得三条推论。证毕。
\end{proof}

\paragraph{采样权重的矩实现：Stieltjes 变换、Hankel 主块与节点谱}
沿用结论 \ref{con:xi-terminal-wronskian-kernel-lagrange-sampling-identity} 的设定，令 \(b_0,\dots,b_\kappa\) 为 \(y_{+}\) 的互异实根，并定义采样权重
\[
c_m:=\frac{y(b_m)}{y_{+}'(b_m)}\qquad(m=0,\dots,\kappa).
\]
则有部分分式恒等式
\begin{equation}\label{eq:xi-terminal-wronskian-stieltjes-transform-minus-y-over-yplus}
\bar m(z)
:=\sum_{m=0}^{\kappa}\frac{c_m}{b_m-z}
\;=\;-\frac{y(z)}{y_{+}(z)}.
\end{equation}
并由 \(z\to\infty\) 的系数比较得到 \(\sum_{m=0}^{\kappa}c_m=1\)（亦即结论 \ref{con:xi-terminal-wronskian-top-coefficient-extraction} 在 \(q=y\) 时的最后一式）。
定义矩序列
\[
\mu_n:=\sum_{m=0}^{\kappa}c_m\,b_m^{n}\qquad(n\ge 0),
\]
并记其 \((\kappa+1)\times(\kappa+1)\) Hankel 主块与一次移位块为
\[
H_{\kappa}:=\bigl(\mu_{i+j}\bigr)_{0\le i,j\le \kappa},
\qquad
H_{\kappa}^{+}:=\bigl(\mu_{i+j+1}\bigr)_{0\le i,j\le \kappa}.
\]

\begin{conclusion}[全主 Hankel 行列式的消去闭式与正权符号律]\label{con:xi-terminal-wronskian-hankel-principal-determinant-resultant}
在上述记号下成立严格恒等式
\[
\boxed{\;
\det H_{\kappa}
=(-1)^{\frac{\kappa(\kappa+1)}{2}}\operatorname{Res}(y,y_{+})\; }.
\]
并且在正权口径 \(c_m>0\) 下，\(H_{\kappa}\) 正定，从而有符号律
\[
(-1)^{\frac{\kappa(\kappa+1)}{2}}\operatorname{Res}(y,y_{+})>0.
\]
\end{conclusion}
\begin{proof}[证明要点]
令 Vandermonde 矩阵 \(V=(b_m^{\,i})_{0\le i\le \kappa,\,0\le m\le \kappa}\)，并记 \(C:=\mathrm{diag}(c_0,\dots,c_\kappa)\)。
由 \(\mu_{i+j}=\sum_{m=0}^{\kappa}c_m b_m^{\,i}b_m^{\,j}\) 得
\(
H_{\kappa}=VCV^{\mathsf T}
\)，
因而
\[
\det H_{\kappa}=(\det V)^2\prod_{m=0}^{\kappa}c_m.
\]
另一方面，由 \(y_{+}\) 首一且零点为 \(\{b_m\}\) 得
\[
\operatorname{Res}(y,y_{+})=\prod_{m=0}^{\kappa}y(b_m),
\qquad
\prod_{m=0}^{\kappa}c_m=\frac{\prod_m y(b_m)}{\prod_m y_{+}'(b_m)}.
\]
又有标准恒等式
\[
(\det V)^2=\mathrm{Disc}(y_{+}),
\qquad
\prod_{m=0}^{\kappa}y_{+}'(b_m)=(-1)^{\frac{\kappa(\kappa+1)}{2}}\mathrm{Disc}(y_{+}),
\]
代入并消去 \(\mathrm{Disc}(y_{+})\) 即得所示闭式。正权情形下 \(H_{\kappa}\) 为 Gram 矩阵，故 \(\det H_{\kappa}>0\)，从而得到符号律。证毕。
\end{proof}

\begin{conclusion}[移位 Hankel 线性笔的节点对角化与 \texorpdfstring{$H_{\kappa}$}{H}-自伴随]\label{con:xi-terminal-wronskian-hankel-shift-diagonalization}
令
\[
A_{\kappa}:=H_{\kappa}^{-1}H_{\kappa}^{+}.
\]
则存在精确相似对角化
\[
\boxed{\;
A_{\kappa}=(V^{\mathsf T})^{-1}\,\mathrm{diag}(b_0,\dots,b_\kappa)\,V^{\mathsf T}\;},
\]
从而特征多项式满足
\[
\chi_{A_{\kappa}}(z)=\prod_{m=0}^{\kappa}(z-b_m)=y_{+}(z).
\]
此外 \(A_{\kappa}\) 满足 \(H_{\kappa}\)-自伴随关系
\[
\boxed{\;
A_{\kappa}^{\mathsf T}H_{\kappa}=H_{\kappa}A_{\kappa}\; }.
\]
\end{conclusion}
\begin{proof}[证明要点]
同理由 \(\mu_{i+j+1}=\sum_m c_m b_m b_m^{\,i}b_m^{\,j}\) 得
\(
H_{\kappa}^{+}=V\,\mathrm{diag}(c_0b_0,\dots,c_\kappa b_\kappa)\,V^{\mathsf T}
\)。
与 \(H_{\kappa}=VCV^{\mathsf T}\) 合并可得
\[
A_{\kappa}=(VCV^{\mathsf T})^{-1}V\,\mathrm{diag}(cb)\,V^{\mathsf T}
=(V^{\mathsf T})^{-1}\mathrm{diag}(b)\,V^{\mathsf T}.
\]
特征多项式在相似变换下不变，从而 \(\chi_{A_{\kappa}}=y_{+}\)。
最后因 \(H_{\kappa},H_{\kappa}^{+}\) 对称，有
\(
A_{\kappa}^{\mathsf T}H_{\kappa}=(H_{\kappa}^{-1}H_{\kappa}^{+})^{\mathsf T}H_{\kappa}=H_{\kappa}^{+}=H_{\kappa}A_{\kappa}
\)。
证毕。
\end{proof}

\begin{conclusion}[节点与权重的双重抽取：双正交本征对与二次型闭式]\label{con:xi-terminal-wronskian-hankel-weights-quadratic-extraction}
令 \(e_m\in\RR^{\kappa+1}\) 为标准基。定义右/左本征向量
\[
r_m:=(V^{\mathsf T})^{-1}e_m,\qquad
\ell_m^{\mathsf T}:=e_m^{\mathsf T}V^{\mathsf T}=(1,b_m,\dots,b_m^{\kappa})\qquad(m=0,\dots,\kappa).
\]
则满足：
\begin{enumerate}
  \item \(A_{\kappa}r_m=b_m r_m\) 且 \(\ell_m^{\mathsf T}A_{\kappa}=b_m \ell_m^{\mathsf T}\)，并且 \(\ell_m^{\mathsf T}r_n=\delta_{mn}\)；
  \item \(\{r_m\}\) 关于双线性型 \(\langle u,v\rangle_H:=u^{\mathsf T}H_{\kappa}v\) 两两正交，且
  \[
  r_m^{\mathsf T}H_{\kappa}r_n=c_m\,\delta_{mn};
  \]
  \item 权重可由 \(H_{\kappa}^{-1}\) 在节点向量上的二次型抽取：
  \[
  \boxed{\;
  c_m^{-1}=\ell_m^{\mathsf T}H_{\kappa}^{-1}\ell_m\; }.
  \]
\end{enumerate}
\end{conclusion}
\begin{proof}[证明要点]
第一条由结论 \ref{con:xi-terminal-wronskian-hankel-shift-diagonalization} 的相似对角化直接推出。
第二条由 \(H_{\kappa}=VCV^{\mathsf T}\) 与 \(r_m=(V^{\mathsf T})^{-1}e_m\) 得
\[
r_m^{\mathsf T}H_{\kappa}r_n
=e_m^{\mathsf T}V^{-1}(VCV^{\mathsf T})(V^{\mathsf T})^{-1}e_n
=e_m^{\mathsf T}Ce_n
=c_m\delta_{mn}.
\]
第三条由 \(H_{\kappa}^{-1}=(V^{\mathsf T})^{-1}C^{-1}V^{-1}\) 与 \(\ell_m=V^{\mathsf T}e_m\) 计算即得。证毕。
\end{proof}

\begin{conclusion}[由 \texorpdfstring{$(H_{\kappa},H_{\kappa}^{+})$}{(H,H+)} 的纯行列式重建 \texorpdfstring{$(y,y_{+})$}{(y,y+)}]\label{con:xi-terminal-wronskian-hankel-determinantal-reconstruction}
令 \(u^{\mathsf T}:=e_0^{\mathsf T}H_{\kappa}=(\mu_0,\dots,\mu_\kappa)\) 且 \(v:=e_0\)。
则有 resolvent 恒等式
\[
\boxed{\;
-\frac{y(z)}{y_{+}(z)}=-u^{\mathsf T}(zI-A_{\kappa})^{-1}v
=-e_0^{\mathsf T}H_{\kappa}\,(zH_{\kappa}-H_{\kappa}^{+})^{-1}H_{\kappa}e_0\; }.
\]
并且 \((y,y_{+})\) 可由确定性行列式直接给出：
\[
\boxed{\;
y_{+}(z)=\frac{\det(zH_{\kappa}-H_{\kappa}^{+})}{\det(H_{\kappa})},
\qquad
y(z)=-\frac{\det\!\begin{pmatrix}
zH_{\kappa}-H_{\kappa}^{+} & H_{\kappa}e_0\\
e_0^{\mathsf T}H_{\kappa} & 0
\end{pmatrix}}{\det(H_{\kappa})}\; }.
\]
\end{conclusion}
\begin{proof}[证明要点]
由 \(A_{\kappa}=H_{\kappa}^{-1}H_{\kappa}^{+}\) 得
\(
H_{\kappa}(zI-A_{\kappa})=zH_{\kappa}-H_{\kappa}^{+}
\)，
从而
\[
\det(zI-A_{\kappa})=\frac{\det(zH_{\kappa}-H_{\kappa}^{+})}{\det(H_{\kappa})}.
\]
再由结论 \ref{con:xi-terminal-wronskian-hankel-shift-diagonalization} 的谱同一性可知 \(\det(zI-A_{\kappa})=y_{+}(z)\)，得到 \(y_{+}\) 的行列式表达。
\par
对 \(-u^{\mathsf T}(zI-A_{\kappa})^{-1}v\) 用 Schur 补恒等式
\[
\det\!\begin{pmatrix} M & v\\ u^{\mathsf T} & 0\end{pmatrix}
=-\det(M)\,u^{\mathsf T}M^{-1}v\qquad(M:=zI-A_{\kappa})
\]
并将块矩阵左乘 \(\mathrm{diag}(H_{\kappa},1)\)（仅引入因子 \(\det(H_{\kappa})\)）即可化为所示的 \((zH_{\kappa}-H_{\kappa}^{+})\) 线性笔行列式比值。
该有理函数与 \eqref{eq:xi-terminal-wronskian-stieltjes-transform-minus-y-over-yplus} 具有同一极点集合 \(\{b_m\}\)、同一留数 \(c_m=y(b_m)/y_{+}'(b_m)\)，且在 \(z\to\infty\) 时同为 \(-\mu_0/z+O(z^{-2})\)，故恒等相等，从而得到 \(y\) 的行列式表达。证毕。
\end{proof}

\begin{conclusion}[Wronskian 完备化下的 \texorpdfstring{$\operatorname{Res}$}{Res}--Hankel--\texorpdfstring{$\mathrm{Disc}$}{Disc} 无符号分解]\label{con:xi-terminal-wronskian-hankel-discriminant-signfree-factorization}
在 Wronskian 完备化 \(A=yy_{+}'-y'y_{+}\) 下，有两条无符号精确分解：
\[
\boxed{\;
\operatorname{Res}(A,y_{+})=\det(H_{\kappa})\,\mathrm{Disc}(y_{+}),
\qquad
\operatorname{Res}(A,y)=\det(H_{\kappa})\,\mathrm{Disc}(y)\; }.
\]
\end{conclusion}
\begin{proof}[证明要点]
结论 \ref{con:xi-terminal-wronskian-resultant-discriminant-factorization} 给出
\[
\operatorname{Res}(A,y_{+})
=(-1)^{\frac{\kappa(\kappa+1)}{2}}\operatorname{Res}(y,y_{+})\,\mathrm{Disc}(y_{+}),
\qquad
\operatorname{Res}(A,y)
=(-1)^{\frac{\kappa(\kappa+1)}{2}}\operatorname{Res}(y,y_{+})\,\mathrm{Disc}(y).
\]
再由结论 \ref{con:xi-terminal-wronskian-hankel-principal-determinant-resultant} 的
\(
\operatorname{Res}(y,y_{+})=(-1)^{\frac{\kappa(\kappa+1)}{2}}\det(H_{\kappa})
\)
代入，符号因子抵消即得所示无符号分解。证毕。
\end{proof}

\endinput
